\chapter{评测与可观测性}
\label{chap:evaluation}

\begin{takeaway}[学习目标]
从任务合同建立评测集，区分结果、过程与系统指标；掌握程序断言、模型评审和人工评审的组合；构造 trace，使线上失败能回流为离线回归样本。
\end{takeaway}

\section{没有评测，就没有可控迭代}

换模型、改 prompt、加工具都可能同时改善一类任务并破坏另一类任务。只凭演示和主观聊天无法发现回归。评测的基本循环是：收集代表任务，冻结输入和标准，运行候选版本，比较指标，人工审查分歧，再决定是否发布。

\section{评测单元是一份任务记录}

每条样本建议包含：

\begin{itemize}
  \item 输入目标与环境 fixture；
  \item 可用工具、权限和预算；
  \item 期望结果或允许结果集合；
  \item 必须满足和绝不能发生的条件；
  \item 评分 rubric、权重和证据；
  \item 难度、场景、风险和来源标签。
\end{itemize}

真实任务常不只有一个字符串答案。可以用状态断言，例如“创建了一份草稿但未发送”“最终报告包含五个来源且日期有效”。

\section{三层指标}

\begin{table}[htbp]
\centering
\begin{tabularx}{\textwidth}{>{\bfseries}p{0.2\textwidth} X X}
\toprule
层级 & 衡量内容 & 指标示例 \\
\midrule
结果 & 任务最后是否有用、正确、安全 & 完成率、事实正确率、违规率 \\
过程 & 走法是否合理、可恢复 & 工具选择率、无效循环、重试次数 \\
系统 & 成本与体验是否可接受 & P50/P95 延迟、token、费用、接管率 \\
\bottomrule
\end{tabularx}
\end{table}

最终结果优先级最高。一个 Agent 即使计划文字漂亮，只要业务状态错误就应判失败。过程指标主要用于诊断和成本优化。

\section{评分器的优先顺序}

\begin{enumerate}
  \item \textbf{确定性断言}：schema、文件、数据库状态、测试、数值和引用 URL。
  \item \textbf{领域规则}：覆盖率、合规词、业务计算、权限策略。
  \item \textbf{模型评审}：语义正确性、完整性、表达质量等难以编码的维度。
  \item \textbf{人工评审}：校准主观维度，处理高风险和评分分歧。
\end{enumerate}

能用程序判断的内容不要交给模型。使用 LLM-as-Judge 时，应提供清晰 rubric 与离散等级，要求引用被评答案中的证据，并用人工标注集测量一致性。

\section{避免评测污染和过拟合}

把样本分为开发集、回归集和保留测试集。开发人员可以看开发集；发布门槛使用稳定回归集；保留集不参与日常调 prompt。定期加入线上新失败，删除重复或不再代表业务的样本。

不要只测“正常问题”。至少覆盖：缺字段、矛盾输入、无答案、权限不足、工具超时、提示注入、重复请求、预算不足和用户中途取消。

\section{成对比较与置信区间}

开放式输出用绝对分数容易漂移。比较两个版本时，可以让评审对同一输入做成对选择，并随机左右顺序。报告胜、平、负以及分场景结果，而不是只报单一均值。

小样本的 2 个百分点提升可能只是噪声。保留逐样本结果，使用 bootstrap 置信区间或至少报告样本数与变化样本。高风险发布还应逐条审查所有退化。

\section{最小评测工具}

\code{examples/03\_eval\_harness.py} 演示 JSON 样本、Agent 运行、确定性评分和聚合报告。一个可维护的评测命令应固定：代码版本、模型版本、prompt 版本、数据集版本、随机参数和环境 fixture。

\begin{codeblock}
python3 examples/03_eval_harness.py
# cases=5 passed=5 pass_rate=100.0%
\end{codeblock}

真实模型存在波动，可对关键样本运行多次，报告 pass@1、最坏结果与方差。不要只挑一次最好结果展示。

\section{Trace 是任务的可解释时间线}

一个 trace 包含 run、step/span 与 event。每个 span 记录开始结束、父节点、组件、输入输出摘要、状态、错误、token、成本和 artifact。跨子 Agent 与工具调用传播 trace ID，才能还原端到端路径。

\begin{warningbox}
可观测性数据常包含 prompt、用户文档、工具参数和模型答案，是高敏感数据。默认脱敏、访问控制、加密与设置保留期；不要为了调试永久保存所有原文。
\end{warningbox}

\section{仪表盘与告警}

仪表盘至少按版本、模型、任务类型、工具和租户切分。关注成功率、失败码、延迟分位数、步骤数、成本、缓存命中、重试和人工接管。告警应指向可行动问题，例如某工具的权限错误率突增，而不是笼统“Agent 质量下降”。

\section{线上失败回流}

建立闭环：线上 trace 被用户点踩或规则判失败；脱敏并最小化；工程师标注根因；生成可复现 fixture；加入回归集；修复后验证；发布并监控。没有这条链路，日志只是一座昂贵的数据仓库。

\begin{practice}[建立 50 条黄金集]
从真实或模拟任务中选 50 条，按正常、边界、故障、安全分层。为每条至少写一个确定性断言。比较两个 prompt 或模型版本，输出总体和分层胜负，并逐条说明所有回归。
\end{practice}

\begin{checkpoint}
\begin{itemize}
  \item \checkbox 评测样本包含环境、权限、预算和禁止条件。
  \item \checkbox 能确定性判断的维度没有交给 LLM 评分。
  \item \checkbox 每次结果可追溯到代码、模型、prompt 和数据集版本。
  \item \checkbox 线上失败可以脱敏后转成可复现回归样本。
\end{itemize}
\end{checkpoint}

